% sec-07-conclusions.tex - Conclusions (full)
\section{Conclusions}

A single repository based LLM prompt, which first writes and then executes a schedule
of sub-prompts through a mermaid to draft to full to final pipeline, can generate
every relevant large Phase 1 document and thereby hasten the entire Phase 1 process.
The acceleration is specific and honest: it compresses the administrative and
preparation bucket, where document white space between phases is lost, and leaves the
clinical-operational and fixed regulatory-review buckets untouched. Applied to the
six highest-value document targets, the method turns iterations that took weeks to
months into iterations that take one to four days.

The method is monitorable and verifiable rather than merely fast. Its 24 figures
ground the prose and reproduce identically in LaTeX; its repository, directory, and
file names match the artifacts on disk; and every file is committed to GitHub in real
time. Because each probable risk is reduced to a low residual by a concrete control,
probable benefit exceeds probable risk for enrolled PDAC patients who cannot wait,
and faster, safer document administration translates, through earlier treatment, into
extended progression-free and overall survival.

This application is credible because the LLM trust it depends on was established
incrementally: the 2024 to 2026 application papers, the National Platform framework
\cite{natplatform}, and the prior single-prompt multi-file builds of the Phase 1 and
Phase 2 daraxonrasib PDAC protocols and the law and bill repositories
\cite{phase1protocol, phase2protocol, kawchak_2026_20196639, lawnarrative,
adoptionguide, lawenactment, hr9510v5, bill01v4} together show that the same
single-prompt approach, now turned on the document-generation problem itself, is a
natural and well-supported next step. Paper v1.0; repository v4.2.0.
